\section{Следствия из теоремы об отделимости.}
\begin{theorem}[Мазур]\label{th:mazur}
    Пусть $(X, \|\cdot\|)$ -- нормированное пространство и $M \subset X$ -- выпукло и замкнуто по норме.
    Тогда $M$ слабо замкнуто.
\end{theorem}
\begin{proof}
    Заметим, что для всякой $x \in X \setminus M$ согласно второй теореме об отделимости существуют $f \in X^*$ и $r \in \R$ такие, что
    \[
        \forall y \in M \hookrightarrow \Re f(x) > r \geq \Re f(y).
    \]
    Рассмотрим $U = V(x, f, \Re f(x) - r)$.
    Тогда $U \cap M = \emptyset$, поскольку для всякой $z \in U$
    \[
        \Re f(x) - r > |f(x) - f(z)| \geq \Re f(x) - \Re f(z) \Rightarrow \Re f(z) > r,
    \]
    поэтому $z \notin M$.
    Следовательно, $x$ -- внутренняя точка (в смысле слабой топологии) $X \setminus M$.
    В силу произвольного выбора $x$ верно, что $X \setminus M$ -- слабо открыто, а значит $M$ слабо замкнуто.
\end{proof}
\begin{corollary}\label{cor:separability-corollaries}
    Если $M$ -- выпуклое подмножество нормированного пространства, то его слабое замыкание совпадает с замыканием по норме.
\end{corollary}
\begin{proof}
    Так как $[M]_{\|\cdot\|}$ -- выпукло (в локально-выпуклых пространствах замыкание не портит выпуклость) и сильно замкнуто, то по теореме Мазура оно и слабо замкнуто.
    Поскольку $M \subset [M]_{\|\cdot\|}$, то $[M]_{w} \subset [M]_{\|\cdot\|}$.
    Но слабо замкнутое множество всегда и по норме замкнуто, следовательно $[M]_{w} = [M]_{\|\cdot\|}$.
\end{proof}
\begin{corollary}
    Пусть $(X, \|\cdot\|)$ -- нормированное пространство и $x_n \xrightarrow{w} x$.
    Тогда существует $\{y_m\} \subset \conv \{x_n\}$ такая, что
    \[
        y_m \xrightarrow{\|\cdot\|} x.
    \]
\end{corollary}
\begin{proof}
    Заметим, что $M = \conv\{x_n\}$ -- выпуклое множество.
    А значит $[M]_{w} = [M]_{\|\cdot\|}$.
    В тоже время $x \in [M]_{w}$, поэтому $x \in [M]_{\|\cdot\|}$.
    Но тогда, очевидно, существует $\{y_m\} \subset M$ сходящееся к $x$ по норме.
\end{proof}
\begin{corollary}
    Пусть $X$ -- рефлексивно, $M \subset X$ -- выпуклое, сильно замкнутое и ограниченное подмножество,
    а $J\colon M \to \R$ -- сильно непрерывный выпуклый функционал.
    Тогда существует $x \in M$ такой, что $J(x) = \inf_{M} J$.
\end{corollary}
\begin{proof}
    По определению инфинума существует $\{x_n\}$ -- минимизирующая последовательность -- такая, что
    \[
        J(x_n) \rightarrow \inf_{M} J = J_* \in [-\infty, +\infty).
    \]
    В силу теорем Банаха-Алаоглу и Эберлейна-Шмульяна $M$ -- слабый секвенциальный компакт, а значит существует слабо сходящаяся к некоторому $x \in [M]_{w}$ подпоследовательность $\{x_{n_k}\}$.
    По теореме Мазура $[M]_{w} = [M]_{\|\cdot\|} = M$.
    В частности, $J(x_{n_k}) \rightarrow J_*$.
    Значит $\forall \gamma > J_*$ существует $K_\gamma$ такое, что $\forall k \geq K_\gamma$ верно
    \[
        J_{*} \leq J(x_{n_k}) < \gamma.
    \]
    Согласно предыдущему следствию существует $\{y_m\} \subset \conv\{x_{n_k}; k \geq K_\gamma\} \subset M$ такая, что $y_m \xrightarrow{\|\cdot\|} x$.
    Поскольку функционал является сильно непрерывным, то $J(y_m) \rightarrow J(x)$.
    Заметим, что $y_m = \sum_{j} t_j x_j$, где $t_j \geq 0$ и $\sum_j t_j = 1$.
    Поэтому в силу выпуклости $J\colon$
    \[
        J(y_m) = J\left(\sum_{j} t_j x_j\right) \leq \sum_{j} t_j J(x_j) < \sum_{j} t_j \gamma = \gamma.
    \]
    А значит и $J(x) \leq \gamma$.
    Но $\gamma$ -- произвольное (большее $J_*$), а значит $J_* = J(x)$.
\end{proof}
\begin{corollary}
    Пусть $X$ -- рефлексивно, $M \subset X$ -- выпуклое и сильно замкнутое множество,
    а $J\colon M \to \R$ -- сильно непрерывный выпуклый функционал такой, что существует $r \in \R\colon$ $L_r = J^{-1}((-\infty, r])$ -- ограничено и непусто.
    Тогда существует $x \in M$ такой, что $J(x) = \inf_{M} J$.
\end{corollary}
\begin{proof}
    Заметим, что $\inf_{M} J = \inf_{L_r}J$.
    Так как $L_r$ -- непусто, то существует некоторое $x \in L_r$ такое, что $J(x) \leq r$.
    Следовательно, $\inf_M J \leq r$.
    Пусть $\{x_m\}$ -- минимизирующая последовательность.
    Тогда $J(x_m) \rightarrow \inf_M J$.
    Начиная с некоторого номера $M$ последовательность $\{J(x_m)\}_{m \geq M}$ будет лежать не правее $r$, а значит лежать в $L_r$.
    Поэтому $\inf_M J \geq \inf_{L_r} J$.
    Но инфинум по меньшему множеству не мог уменьшится, а значит $\inf_M J \leq \inf_{L_r} J$, что и даёт нужное равенство.

    Теперь мы можем применить предыдущее следствие к $L_r$ и $J$.
\end{proof}
\begin{corollary}\label{cor:metric-projection}
    Пусть $X$ -- рефлексивное, $M$ -- непустое выпуклое и сильно замкнутое подмножество $X$.
    Тогда для всякого $z \in X$ существует элемент наилучшего приближения в $M$, то есть на некотором $y \in M$ достигается $\rho(z, M)\colon$
    \[
        \rho(z, M) = \|z - y\|.
    \]
\end{corollary}
\begin{proof}
    Рассмотрим $J(x) = \|x - z\|$.
    Там как $M \neq \emptyset$, то в нём лежит некоторый $y_0$.
    Положим $r = \|z - y_0\|$.
    Очевидно, что $L_r$ ограничен по норме и непуст (как минимум в нём лежит $y_0$).
    Тогда мы можем применить предыдущее следствие и получить требуемое утверждение.
\end{proof}
\begin{definition}\label{def:annulators}
    Пусть $X$ -- нормированное пространство, и $L \subset X, N \subset X^*$ -- подпространства.
    \begin{itemize}
        \item Правым аннулятором $L$ назовём $L^\perp = \{f \in X^* \mid f(x) = 0 \ \forall x \in L\}$.
        \item Левым аннулятором $N$ назовём $^\perp N = \{x \in X \mid f(x) = 0 \ \forall f \in N\}$.
    \end{itemize}
\end{definition}
\begin{note}
    Заметим, что $L^\perp$ -- это подпространство $X^*$, замкнутое по $\tau_{*w}$.
    Это действительно так, поскольку
    \[
        L^{\perp} = \bigcap_{x \in L} J_x^{-1}(\{0\}),
    \]
    где $J\colon X \to X^{**}$ -- каноническая изометрия Банаха, а $J_x$ -- *-слабо непрерывный функционал на $X^*$ (по определению).

    Заметим, что $^\perp N$ -- сильно замкнутое подпространство $X$ (а, значит, и слабо замкнутое).
    Аналогично, очевидно, что
    \[
        ^{\perp} N = \bigcap_{f \in N} f^{-1}(\{0\}),
    \]
    а всякий $f \in N \subset X^*$ сильно непрерывен, поэтому $^\perp N$ сильно замкнут как пересечение сильно замкнутых (а так как это -- подпространство, то, по теореме Мазура, оно ещё и слабо замкнуто).
\end{note}
\begin{corollary}\label{cor:annulator_equations}
    Верно следующее:
    \begin{itemize}
        \item $^\perp(L^\perp) = [L]_{\|\cdot\|} = [L]_{w}$.
        \item $(^\perp N)^\perp = [N]_{*w} \supset [N]_{\|\cdot\|}$.
    \end{itemize}
\end{corollary}
\begin{proof}
    Покажем первое равенство.
    Очевидно, что $L \subset ^{\perp}(L^\perp)$.
    А значит, в силу сильной замкнутости $^\perp (L^\perp)$, и $[L]_{\|\cdot\|} \subset ^\perp (L^\perp)$.

    Пусть $x \in ^\perp(L^\perp) \setminus [L]_{\|\cdot\|}$.
    Тогда в силу второй теоремы об отделимости существуют $f \in X^*$ и $r \in \R$ такие, что
    \[
        \forall y \in [L]_{\|\cdot\|} \hookrightarrow \Re f(x) > r \geq \Re f(y).
    \]
    Следовательно, $\Re f$ и $\Im f$ ограничены сверху $r$ на $L$.
    Если существует $y_0 \in L$ такое, что $f(y_0) \neq 0$, то мы приходим к противоречию.
    А значит $f|_L = 0$, и, поэтому, $f \in L^\perp$.следовательно,
    Но $f(x) > r \geq 0$ -- противоречие с тем, что $x \in ^\perp(L^\perp)$.

    Покажем второе равенство.
    Аналогично, $N \subset (^\perp N)^\perp$, а значит и $[N]_{*w} \subset (^\perp N)^\perp$.
    Более того, верна следующая цепочка включений:
    \[
        [N]_{\|\cdot\|} \subset [N]_{*w} \subset (^\perp N)^{\perp}.
    \]
    Заметим, что так как $\tau_{*w}$ -- локально-выпуклая топология на $X^*$, то мы можем использовать теорему об отделимости.
    Пусть $\exists f \in (^\perp N)^\perp \setminus [N]_{*w}$.
    По теореме об отделимости $\exists \psi \in (X^*, \tau_{*w})^* = \Phi(X)$ и $r \in \R$ такие, что
    \[
        \forall h \in N \hookrightarrow \psi(f) > r \geq \psi(h).
    \]
    Действуя аналогично первому случаю получаем то, что $\psi|_{N} = 0$, и тогда, так как $\psi = \Phi_x$ для некоторого $x \in X$ (под действием изометрии Банаха $\Phi$), то $x \in ^\perp N$ -- противоречие.
\end{proof}
\begin{notet}[В ТВП замыкание выпуклого множества -- выпукло]
    Пусть $E$ — ТВП и $C$ — выпукло. Рассмотрим $x, y \in [C]$ и $t \in (0, 1)$. Пусть $z = tx + (1 - t)y$. Для того, чтобы показать, что $z$ лежит в замыкании $C$ достаточно показать, что всякая окрестность $z$ имеет непустое пересечение с $C$.

    Пусть $U$ — окрестность $z$. Рассмотрим $A_t(u, v) = tu + (1 - t)v$. Заметим, что $A_t$ — непрерывно, так как умножение на скаляр и сложение в ТВП непрерывны. В частности, $A_t$ — непрерывно в точке $(x, y)$ и $z = A_t(x, y)$. Тогда у нас существуют окрестности $V$ и $W$ точек $x$ и $y$ соответственно такие, что $A_t(V, W) \subset U$. Так как $x, y \in [C]$, то $V \cap C$ и $W \cap C$ — непусты. Пусть $a$ лежит в $V \cap C$ и $b$ лежит в $W \cap C$. Тогда $A_t(a, b) \in U$.
    
    Но так как $C$ — выпукло, то $A_t(a, b) = at + (1 - t)b \in C$, а значит $U \cap C$ — непусто. Так как $U$ была выбрана произвольно, то $z \in [C]$ и $[C]$ — выпукло,
\end{notet}
\begin{note}[Теорема Голдстайна]\label{th:goldstein}
    Пусть $(X, \|\cdot\|)$ -- нормированное пространство, $\Phi\colon X \to X^{**}$ -- каноническая изометрия Банаха.
    Тогда $[\Phi(X)]_{*w} = X^{**}$

    Покажем это.
    Так как
    \[
        [\Phi(X)]_{*w} = (^\perp \Phi(X))^{\perp},
    \]
    а $^\perp \Phi(X) = \{f \in X^* \mid \Phi_x f = 0 \forall x \in X \} = \{f \in X^* \mid f(x) = 0 \forall x \in X\} = \{0\}$,
    то
    \[
        [\Phi(X)]_{*w} = \{0\}^{\perp} = X^{**}.
    \]
\end{note}
\begin{note}\label{cor:reflectia-criterion}
    Заметим, что если $X$ -- банахово, то в силу изометричности $\Phi$ множество $\Phi(X)$ -- сильно замкнуто в $X^{**}$.
    При этом, если $X$ -- нерефлексивно, то $\Phi(X)$ будет *-слабо в $X^{**}$ не замкнуто (в силу теоремы Голдстайна).
\end{note}
\begin{corollary}
    Пусть $X$ -- нормированное пространство такое, что $B_1[0] \subset X$ -- слабый компакт.
    Тогда $X$ -- рефлексивно.
\end{corollary}
\begin{proof}
    Известно, что $\Phi\colon (X, \tau_w) \to (\Phi(X), \tau_{*w})$ является топологически непрерывным.
    А значит $\Phi B_1[0]$ является *-слабым компактом в $(\Phi(X), \tau_{*w})$, то есть и в $(X^{**}, \tau_{*w})$.
    В тоже время $\Phi B_1[0]$ -- выпукло.
    Поскольку *-слабая топология хаусдорфова, то $\Phi B_1[0]$ -- *-слабо замкнуто в $X^{**}$.

    Пусть $\phi \in B_1^{**}[0] \setminus \Phi B_1[0]$.
    Тогда мы можем воспользоваться теоремой об отделимости и получить $\xi \in (X^{**}, \tau_{*w})^* = X^*$, а значит $\xi(\psi) = \psi(f)$ для некоторого $f \in X^*$.
    Тогда:
    \[
        \forall x \in B_1[0] \hookrightarrow \Re \xi(\psi) > r \geq \Re \xi(\Phi_x).
    \]
    Перепишем в более удобном виде:
    \[
        \forall x \in B_1[0] \hookrightarrow \Re \psi(f) > r \geq \Re \xi(\Phi_x).
    \]
    Возьмём супремум по $x \in B_1[0]$ и получим, что
    \[
        \|\psi \| \|f\| > r \geq \|\Re f\| = \|f\| \Rightarrow \|\psi \| > 1.
    \]
    Пришли к противоречию.
\end{proof}

    \begin{theorem}\label{th:refref}
        Пусть $X$ -- банахово пространство.
        Тогда следующие утверждения эквивалентны:
        \begin{itemize}
            \item $X$ -- рефлексивно.
            \item $X^*$ -- рефлексивно.
        \end{itemize}
    \end{theorem}
    \begin{proof}
        Пусть $J_X\colon X \to X^{**}$ -- каноническая изомтерия Банаха и $X^*$ -- рефлексивно.
        Заметим, что $J_X(X)$ -- замкнутое подпространство $X^{**}$.
        По определению, $X$ -- рефлексивно, если $J_X(X) = X^{**}$.
        Для замкнутого подпространства это эквивалентно тому, что $J_X(X)^\perp = 0 \subset X^{***}$.

        Пусть $\Phi \in X^{***}$. Если $\Phi \in J_X(X)^\perp$, то $\Phi(J_X x) = 0$ при всяком $x \in X$. Так как $X^*$ -- рефлексивно, то $\Phi = J_{X^*}f$ для некоторого $f \in X^*$.
        Поэтому
        \[
            0 = \Phi(J_X x) = (J_{X^*} f)(J_X x) = (J_X x)(f) = f(x).
        \]
        Так как это верно при всех $x \in X$, то $f = 0$, а значит и $\Phi = 0$ -- и $J_X(X)^\perp = 0$.

        Пусть $X$ -- рефлексивно. Положим $f = J_X^* \Phi$ при всяком $\Phi \in X^{***}$.
        Заметим, что
        \[
            J_{X^*}(f)(J_X x) = J_{X^*}(J_{X}^*\Phi)(J_X x) = (J_X x)(J_{X}^{*}\Phi) = (J_{X}^{*}\Phi)(x) = \Phi(J_X x),
        \]
        что и даёт нам искомую сюръективность в силу того, что $X^{**} = J_X X$.
    \end{proof}
